\chapter{Imports, Automated Origin and Negative Results}\label{ch:imports}

Large systems are built partly by machines: bulk imports, extraction pipelines and generated contributions. They raise three problems for a claim graph. Volume can look like consensus. Machine production can pass for independent human confirmation. And what is measured, imported or generated tends to be positive findings, because those are what sources publish. This chapter states how each is treated, and proves the properties that follow from the model.

\section{Origin modes and batches}\label{sec:origin}

Every act record declares an origin mode: \textsf{human}, \textsf{automated}, \textsf{assisted} or \textsf{imported}. The declaration is part of the trace (Definition~\ref{def:trace}), and a record without one is trace-null and quarantined. The mode is displayed and can be filtered in queries. It does not enter the labeling (Proposition~\ref{prop:separation}).

\begin{definition}[Batch]\label{def:batch}
A \emph{batch record} references a set $M$ of records, states the procedure that produced them, the source or sources it drew on, and the origin mode, and asserts a \emph{common-origin claim}: that the members of $M$ lie in one independence family. A record is a \emph{member} of the batch if the batch references it.
\end{definition}

The common-origin claim is an equivalence claim in the sense of Chapter~\ref{ch:distinction}. It is an ordinary claim, and it is open to objection, but the burden lies in the direction the design intends: members of a batch are one family unless a distinction claim with grounds says otherwise.

\begin{designrule}[Batches are one family by default]\label{rule:batchfamily}
The default family assignment for members of a batch is a single family. Assigning members to distinct families requires a distinction claim stating why their origins are independent, filed as an ordinary act and open to objection. Records of different batches are in different families only if the procedures or sources are declared to differ.
\end{designrule}

\begin{proposition}[An import is one family]\label{prop:importfamily}
Let a batch have $n$ members, each of which is an evidence item, and suppose the common-origin claim of the batch stands. Then the members lie in one independence family, and, for any claim, their joint contribution to the replication coordinate $r^+$ or $r^-$ of the status vector is at most $1$, whatever $n$ is.
\end{proposition}

\begin{proof}
By definition the family identifier is the class under the equivalence generated by equivalence claims that stand, and the common-origin claim identifies all members. The coordinates $r^\pm$ count independence families and not records (Definition~\ref{def:status}), and Proposition~\ref{prop:indep} shows that further members of an existing family do not raise the count.
\end{proof}

The proposition states, for a batch, what Chapter~\ref{ch:traces} states for automated origin generally: independence, replication and contributor count are never inferred from record count.

\section{Reversible activation}\label{sec:reversible}

An import is activated on a surface by being admitted there. Its withdrawal is a set of withdrawals of its arguments, and this is as clean as a single withdrawal.

\begin{corollary}[Batch withdrawal is deletion up to labeling]\label{cor:batchwithdraw}
Let $a_1,\dots,a_k$ be arguments and let $\ledger'$ extend $\ledger$ by a withdrawal of each. For $x\in\sigma(a_i)$ let $Z^{(i)}_x$ be the set of Theorem~\ref{thm:withdraw} for $a_i$, and let $Z_x=\bigcup_{i:\,x\in\sigma(a_i)}Z^{(i)}_x$. Then every member of $Z_x$ is $\OUT$ at $x$ in $\ledger'$, and every other argument has at $x$ the label it has in the framework obtained from $\mathcal A_x(\ledger)$ by deleting $Z_x$ and all defeats involving its members. At units in no $\sigma(a_i)$ no label changes.
\end{corollary}

\begin{proof}
The proof of Theorem~\ref{thm:withdraw} applies verbatim with $Z_x$ in place of the single set. Each $w_{a_i}$ has no defeaters, so it is $\IN$ and defeats $a_i$; hence every $a_i$ is $\OUT$ at units of $\sigma(a_i)$, and every argument having some $a_i$ as a sub-argument is $\OUT$ there by Lemma~\ref{lem:hereditary}(ii). All members of $Z_x$ are therefore $\OUT$. Delete them one at a time; by Lemma~\ref{lem:outinert} no other label changes and the remaining members stay $\OUT$. The withdrawal nodes then defeat nothing that remains and are isolated $\IN$ nodes.
\end{proof}

A defective import can therefore be withdrawn as a unit, and the graph afterwards is what it would have been had the import never entered, at the level of labels. The record of the import and of its withdrawal stays in the ledger, along with everything anyone did in the interim.

\section{Negative and failed work}\label{sec:negative}

Knowledge bases favor positive assertions because sources do. The specification makes negative and failed work first-class by giving each kind an object in the model.

\begin{center}
\small
\begin{tabularx}{\textwidth}{@{}l X@{}}
\toprule
kind of negative work & object in the model\\
\midrule
failed replication & \textsf{FAILS\_TO\_REPLICATE} from a declared attempt whose outcome does not reproduce the specified result (Chapter~\ref{ch:relations}); counted in $r^-$; supports no opposing kernel\\
null detection & an evidence item whose kernel states that a test under its design did not detect the effect; opposes nothing by itself\\
bounded effect & an evidence item with kernel $\nk\kap$, ``the effect is below the declared bound'', from an equivalence or noninferiority design with adequate precision\\
opposite effect & an evidence item supporting a contrary directional kernel; reaches $\nk\kap$ through an entailment warrant\\
withdrawn argument & a withdrawal record; the argument stays addressable (Theorem~\ref{thm:withdraw})\\
abandoned operationalization & a revision superseding the kernel, with the earlier version preserved (Theorem~\ref{thm:lineage})\\
defeated warrant & an undercutting argument that stands, with its region\\
unsuccessful search & an evidence item reporting the protocol, corpus and date, with kernel stating that no match was found; an inference to a claim of absence needs a warrant of search adequacy, with backing, open to undercutting\\
expected evidence not found & the same, with the expectation stated as a claim\\
\bottomrule
\end{tabularx}
\end{center}

Four things are called a negative result and they are not the same thing. Let $\kap$ be a kernel of the form ``the effect is at least $\delta$'', with the bound $\delta$ part of its operational meaning, and let $\mathsf{det}(\kap,D)$ be the kernel ``a test of $\kap$ under design $D$ detects it''.

\begin{definition}[Negative outcomes]\label{def:negoutcomes}
A study of design $D$ over a scope $\sigma_e$ reports one of the following, each as an evidence item with the kernel shown.
\begin{enumerate}[label=(\roman*)]
\item \emph{Null detection}: the kernel $\nk{\mathsf{det}(\kap,D)}$. The test did not reject its null under the design.
\item \emph{Bounded effect}: the kernel $\nk\kap$. The reported interval excludes effects of size $\delta$ or more, which is a claim the design warrants only if it is an equivalence or noninferiority design with adequate precision at that bound.
\item \emph{Opposite effect}: a kernel $\kap^{\downarrow}$ stating an effect of the opposite sign of size at least $\delta$.
\item \emph{Failed replication}: an evidence item that declares an attempt to follow the protocol of an earlier item $e_0$ and whose outcome kernel is not that of $e_0$. It records that the specified result was not reproduced.
\end{enumerate}
\end{definition}

\begin{proposition}[What each negative outcome can support]\label{prop:negsupport}
\begin{enumerate}[label=(\alph*)]
\item A null detection item has a kernel that is neither $\kap$ nor $\nk\kap$. It is therefore not the conclusion kernel of any argument for $\nk\kap$, and no attack arises between it and any argument for $\kap$. An argument from it to $\nk\kap$ exists only through an ordinary warrant, for instance that a design of adequate power would have detected an effect of size $\delta$ if present, with a backing that states the power, and that warrant is open to undercutting.
\item A bounded-effect item supports arguments for $\nk\kap$ directly, and an opposite-effect item supports them through an entailment warrant, valid in all models in which $\kap^{\downarrow}$ has its stated meaning, from $\kap^{\downarrow}$ to $\nk\kap$.
\item A failed replication contributes to the replication coordinate $r^-$ and to nothing else; it is not evidence for any kernel.
\end{enumerate}
\end{proposition}

\begin{proof}
(a) Attack is defined from equality and negation of kernels (Definition~\ref{def:attack}), and $\nk{\mathsf{det}(\kap,D)}$ is neither equal to nor the negation of $\kap$. The warrant is an ordinary warrant with premise kernel $\nk{\mathsf{det}(\kap,D)}$ and conclusion $\nk\kap$ and is subject to (W1)--(W5). (b) The first is by kernel equality. For the second, if $x\in[\kap^{\downarrow}]_M$ then the effect at $x$ is of the opposite sign, hence below $\delta$, so $x\in[\nk\kap]_M$ in every model with that meaning. (c) By the definition of $r^-$ (Definition~\ref{def:status}) and of \textsf{FAILS\_TO\_REPLICATE}.
\end{proof}

The absence of a positivity bias in the semantics is a precise property.

\begin{proposition}[Evaluation is symmetric under exchange of a kernel and its negation]\label{prop:nopositivity}
Let $\iota$ be a bijection on kernels that commutes with negation, that is, $\iota(\nk\kap)=\nk{\iota(\kap)}$. Extend $\iota$ to a ledger $\ledger$ by renaming the kernels of all claims, evidence items and warrants. Then the argument frameworks of $\ledger$ and $\iota\ledger$ are isomorphic under the renaming, and $\Lab(\iota a,x)=\Lab(a,x)$ for every argument $a$ and unit $x$, when the policy is renamed correspondingly.
\end{proposition}

\begin{proof}
Attack (Definition~\ref{def:attack}) is defined from scopes, from the structure of argument trees, and from whether two kernels are equal or are negations of each other. All three are preserved by $\iota$, which is a bijection commuting with negation and leaving scopes untouched. Denial kernels of targets are renamed consistently. Preferences and admissibility refer to arguments and evidence kinds and are transported by the renaming. The renamed framework is then isomorphic, and the grounded labeling, being defined from the graph alone, is preserved.
\end{proof}

In particular, taking $\iota(\kap)=\nk\kap$ on all kernels exchanges every finding with its opposite. A bounded-effect finding is exactly as strong, in the semantics, as a positive one of the same scope, kind and independence. A null detection is not the mirror image of a positive detection; it is a different kernel, and Proposition~\ref{prop:negsupport}(a) says what it takes to use it. Any asymmetry a dossier shows comes from the records, that is, from what has been inscribed, and none from the rules.

\begin{example}[An unsuccessful search]\label{ex:search}
An evidence item reports a search: protocol $P$, corpus $C$, date $T$, and finds no item matching the query. Its kernel states that no match exists in $C$ at $T$, over the scope of units $C$ covers. To conclude that a phenomenon is absent from the units $C$ represents, an argument needs a warrant of search adequacy, that a search of this kind would have found the phenomenon if present, with a backing that names the search's sensitivity. That warrant can be undercut, and its scope is $\sigma_w$. The dossier thereby records both the search and the assumption that turns it into a finding of absence, and a reader can see which is contested.
\end{example}

A failed attempt does not by that fact lower an actor's standing. A well-executed failure is among the most valuable acts in the system, and the credit dimensions below treat it so.

\section{Credit as a vector}\label{sec:credit}

Credit for argumentative acts is not a scalar, for the reason status is not (Theorem~\ref{thm:noscalar}). Each act $a$ receives a \emph{credit vector}
\[
\rho(a)=(q,\ \delta,\ r,\ v,\ s),
\]
whose components are evidential quality (how well the act's leaves are anchored and how strong the rung they support), decisiveness (how much of the label of the target the act changes), repair value (how much of a defeated claim the act saves, as a survivor), verification (survival under independent challenge or confirmation), and scope precision (agreement between the scope the act asserts and the scope that ultimately stands). Credit attaches \emph{after} downstream events: an objection may earn submission credit on inscription and verification credit only once it survives challenge.

\begin{definition}[Independent validating event]\label{def:validating}
A validating event for an act $a$ is a confirmation of it, a failed attempt to defeat it, or a replication of its finding. It is \emph{independent} if the record that constitutes it lies in an independence family different from that of $a$.
\end{definition}

\begin{proposition}[No self-validation]\label{prop:noself}
The verification component $v(a)$ of the credit of $a$ counts only independent validating events. Consequently adding any number of records in the family of $a$, whether by the author of $a$ or by any other actor, leaves $v(a)$ unchanged.
\end{proposition}

\begin{proof}
By Definition~\ref{def:validating} an event in the family of $a$ is not independent and is not counted.
\end{proof}

The guarantee is only as strong as the family structure that supports it. Family formation is by equivalence claims, which are themselves arguable, so that a ring of coordinated actors that has not been shown to depend on one another counts as independent until someone files and sustains the equivalence claim. The specification does not pretend otherwise, and Chapter~\ref{ch:gates} lists the audits that make such a claim easier to file.

Credit is used for capability decisions and for the allocation of attention. It is derived from labels, and it enters no label (Proposition~\ref{prop:oneway}).

\begin{remark}[Redundant acts]\label{rem:redundant}
It is natural to propose that credit be proportional to the informational change an act makes to the accepted graph, with diminishing credit for structurally redundant acts. The specification adopts this as a governance heuristic and does not define the measure. No theorem in this monograph depends on it, and none should until an information measure on argument graphs is defined that supports one.
\end{remark}

\section{Attention is the scarce resource}\label{sec:attention}

Spam is damaging when it consumes other people's attention, so the scarce thing to charge for is promotion into shared attention surfaces, not speech. The design is unbounded inscription with bounded attention, in place of bounded inscription with unaccountable discretion.

\begin{designrule}[Attention controls]\label{rule:attention}
A surface may combine the following controls, each stated in its policy and each visible in dispositions and reasons: review bonds returned when a contribution survives screening; sponsorship by an actor holding a capability for the relevant family; indexing bandwidth earned progressively; duplicate detection by content hash; bundling of repetitive submissions into a batch; random audits; and reductions of credit for acts that systematically waste review effort. None of them gates inscription (Proposition~\ref{prop:openinscription}), and none changes a label (Proposition~\ref{prop:capnostanding}).
\end{designrule}
